% -*-coding: utf-8 -*-

\defaultfont

\BiAppendixChapter{结~~~~论}{Conclusions}



针对微小型水下运载器平台空间受限、承载能力与抗流场干扰能力有限等问题，
本文提出了一种新型水下绳驱机械臂设计方案。该方案突破了传统刚性机械臂自重偏大、
对微小型水下运载器平台影响大的设计局限，以绳索传输动力，将自重较大的驱动电机后置，减少了移动质量。
全文以系统设计、理论建模与试验验证展开，主要研究工作与成果总结如下：

\begin{enumerate}
    \item \textbf{提出并实现了一种高负载自重比的水下绳驱机械臂总体结构方案}，
    研制了一个高负载自重比的绳驱机械臂物理样机。在构型布局上，
    将电机及主控板统一布置于基座后端的独立密封舱内，
    减小了运动部件质量。针对运载器底盘腹部的狭小安装空间，
    开发了纯滚动双铰链机构以实现大臂的完全折叠
    设计了滑轮组机构与单轴差动缠绕导向方案，
    提升了承载能力。
    通过理论力学计算与 ANSYS 有限元分析，
    验证了核心滑轮轴、连接件及电机驱动轴在极限工况下的结构强度满足水下作业的安全需求。
    
    \item \textbf{推导了耦合约束下的运动学机理与轨迹平滑算法}，
    运用改进D-H法推导正运动学方程，并借助数值迭代法解决了多关节联动逆解难题。
    为了机械臂工作平稳，对比了梯形速度规划与五次多项式算法的插补效果
    并在关节空间与笛卡尔空间分别实现了平滑的直线及圆弧运动轨迹生成。
    MATLAB 仿真结果表明，所提算法有效抑制了速度和加速度突变，
    显著减小了机械臂运动对水体的冲击扰动。
    
    \item \textbf{构建了基于单目视觉的目标定位与自主抓取策略}，
    建立了水下环境的单目相机模型，
    基于 ArUco 标签和“眼在手外”构型，
    通过 Tsai-Lenz 两步法完成了高精度的手眼标定。
    求得了图像像素空间到机械臂物理空间的坐标变换矩阵，
    并设计了基于 ROS 2 架构的抓取任务节点，
    实现了对目标物体的动态定位与抓取动作的自主规划。
    
    \item \textbf{实施了整机软硬件联调与闭环测试}，
    构建了“树莓派+STM32”的主从式计算架构，
    并在底层部署了抗干扰校验的通信架构。
    在视觉抓取及定轨迹运行实验中
    机械臂具备良好的运动响应性能与定位精度，
    视觉引导下的抓取动作连贯，
    验证了整套软硬件系统的可行性与理论模型的正确性。
\end{enumerate}

尽管本文在水下绳驱机械臂的设计与控制上取得了一定成果，
本文的工作仍存在一定的局限性。为进一步推动水下绳驱机械臂的工程化与实用化，
未来可在以下几个方面开展深入研究：

\begin{enumerate}
    \item \textbf{绳索柔性动力学建模与误差补偿}：
    当前的运动学建模主要基于理想刚性假设，
    未充分考虑绳索在长期变载荷工况下的弹性形变与蠕变效应。
    后续研究可引入柔性体动力学模型，并结合张力传感器的实时数据，
    设计动态误差补偿算法，从而进一步突破机械臂的绝对定位精度瓶颈。
    
    \item \textbf{基于连续视觉伺服的动态抓取}：
    现有的视觉引导抓取方案属于单次静态定位，
    即相机单次解算目标坐标后直接下发指令，
    运动过程中缺乏对末端与目标相对位姿的持续监测。
    未来应引入基于图像或基于位置的闭环视觉伺服机制，
    以实现对水下动态目标的实时跟踪与高精度抓取。
    
    % \item \textbf{水动力学参数辨识与抗扰动控制}：
    % 在真实水下作业时，机械臂的各连杆不可避免地会受到流体阻力、
    % 附加质量力等高度非线性因素的影响。未来的底层算法需结合计算流体力学
    % 进行水动力学参数辨识，
    % 并开发滑模控制或自适应鲁棒控制等高级算法，以抵抗复杂水流扰动，
    % 保障轨迹跟踪的平稳性。
    
    \item \textbf{真实水下环境的样机海试与系统联调}：
    目前样机的综合性能验证主要在陆上常压实验室内完成，
    尚未全面考核其动静密封可靠性与涉水工作寿命。
    下一步的研究计划是将该机械臂搭载于真实的微小型水下运载器
    平台上，开展水池及开放水域的实地联调测试。
\end{enumerate}